\documentclass[10pt,a4paper]{amsart}
\usepackage[margin=19mm]{geometry}
\usepackage{amsmath,amssymb,mathtools}
\usepackage[scr=rsfs,cal=euler]{mathalfa}
\usepackage{xfrac}
\usepackage[unicode,hidelinks,bookmarks=false]{hyperref}
\newcommand{\ie}{i.e.\ }
\newcommand{\cf}{cf.\ }
\newcommand{\resp}{respectively\ }
\newcommand{\Subsection}{\subsection}
\providecommand{\nobreakdash}{\nobreak\hbox{-}}
\setlength{\emergencystretch}{2em}
\multlinegap0pt
\usepackage{bm}
\usepackage{bbm}
\usepackage{dsfont}
\usepackage{xspace}
\usepackage{cancel}
\usepackage{mathtools}
\usepackage[shortlabels]{enumitem}
\usepackage{amsthm}
\makeatletter
\@ifundefined{claim}{\newtheorem{claim}{Claim}}{}
\makeatother
\makeatletter
\newcommand{\pbar}{\overline{\varphi}}
\expandafter\ifx\csname proofc\endcsname\relax
\newenvironment{proofc}{\begin{proof}[Proof of Claim]}{\end{proof}}
\fi
\makeatother
\title[Total coloring graphs with large minimum degree]{Total coloring graphs with large minimum degree}
\author{Owen Henderschedt, Jessica McDonald, Songling Shan}
\date{Source: arXiv:2507.05548v1 (CC BY 4.0)}
\begin{document}
\maketitle

\noindent\emph{Source and licence.} Total coloring graphs with large minimum degree. Version arXiv:2507.05548v1 (CC BY 4.0), distributed under the Creative Commons Attribution 4.0 International licence, as recorded in the arXiv licence metadata for this exact version and shown on the abstract page https://arxiv.org/abs/2507.05548v1. Full rights notice: licenses/arxiv-2507.05548-CC-BY-4.0.txt.\par\smallskip

\noindent\textit{Editorial scope.} Counted unit: the Claim whose source label is \texttt{gamma} in the article (source0.tex lines 358--363, source label \texttt{gamma}), delivered together with the article's own complete original proof of it (source0.tex lines 365--437, one \texttt{proof} environment of 73 lines). No mathematics is written, repaired or re-derived: statement and proof are the article's own text line for line. Required context retained uncounted: uwi (lines 343--344). Every definition, display and label of those lines is retained unchanged. Omitted from this package and not redistributed: 1--342, 345--357, 364--364, 438--1375. Nothing inside a retained excerpt is omitted or altered: every retained body is delivered exactly as the article's lines are written, so no omission receipt of any kind accompanies this package. Citations: the retained excerpts carry no citation command at all, so no bibliography entry of the article is transcribed and no reference is redistributed. Reference adaptation: none was needed and none was made; every reference inside the delivered unit resolves inside it, so no unresolved reference or citation is produced. Printed slips kept literal: no printed word, symbol, hypothesis or number of the counted statement or of its proof was altered. Layout and class: the article's own class is \texttt{article} with packages including amsmath, amssymb, amsthm, geometry, enumerate, cases, caption, graphicx, pgf, tikz, pgfplots, tikz-3dplot, fontenc, xcolor; the wrapper is the lane's pinned \texttt{amsart} editorial wrapper at 10pt on A4, which restates the theorem-like environments the retained text needs and adds the article's own macro definitions that the retained text needs (\texttt{\textbackslash pbar}), with extra wrapper packages (bm, bbm, dsfont, xspace, cancel, mathtools, enumitem). The delivered numbering is the unit's own. Assets: no retained span contains a figure environment, a graphics-inclusion command, a PDF-page inclusion, a table or a TikZ picture; no image file is copied into this package, the package declares no asset dependency and no image file is redistributed with it. Licence: version arXiv:2507.05548v1 of this article is distributed under the Creative Commons Attribution 4.0 International licence, as recorded in the arXiv licence metadata for this exact version and shown on the abstract page https://arxiv.org/abs/2507.05548v1. A full rights notice is delivered at licenses/arxiv-2507.05548-CC-BY-4.0.txt. Characters: every retained excerpt is pure ASCII, so the delivered engine, XeLaTeX with its default Latin Modern text font, reports no missing character.
\begin{claim}\label{uwi} Suppose $F_2\subseteq F_1$ corresponds to the union of any number of linear sequences of $F$. Then $\pbar(u)\cap \pbar(w_i) =\emptyset$ for any $w_i\in V(F_2)$.
\end{claim}

\begin{claim}\label{gamma} For any color $\gamma \in [1,k]$,  we have the following statements hold. 
\begin{enumerate}[(1)]
    \item In statement (a),  there are at most  $3$  vertices from $\{w_0, \ldots, w_t\}$ that are all missing $\gamma$ under $\varphi$.
    \item In statement (b),  there is at  most  $1$  vertex from $\{w_0, \ldots, w_t\}$ that is  missing $\gamma$ under $\varphi$.
\end{enumerate}
\end{claim} 

\begin{proofc}  
(1) Suppose for a contradiction that such a color $\gamma$ exists.  Let  $y_1, y_2, y_3,  y_{4}$ be $4$  distinct vertices of $\{w_0, \ldots, w_t\}$ all of which are missing $\gamma$ under $\varphi$. Suppose, without loss of generality, that $y_1$ is the first vertex 
in the order $w_0, w_1, \ldots, w_t$ for which $\gamma\in \pbar(y_1)$.  In particular, this guarantees that $F_1(u, w_0:y_1)$ does not contain any edge that is colored with $\gamma$. 
Let $\alpha\in \pbar(u)$; we know that $\alpha\ne \gamma$ by Claim \ref{uwi}. 



Consider the set $\mathcal{P}$ of maximal $(\alpha,\gamma)$-alternating
path containing at least one vertex from $\{u,y_1,...,y_4\}$. Since $\alpha\in \pbar(u)$ and $\gamma\in \cap_{i=1}^4\pbar(y_i)$, each component of $\mathcal{P}$ is a path where every vertex of $\{u,y_1,...,y_4\}$ is an endpoint of a path. Thus, there are at least $3$ paths in $\mathcal{P}$. By (P1), there is at most one edge of $J_0$ colored $\alpha$ and at most one edge of $J_0$ colored $\gamma$. Hence there must exist a path $P\in \mathcal{P}$ which does not contain any edges of $J_0$. We let $\varphi' = \varphi/P$ and consider two different cases depending on whether $u$ is an endpoint of $P$.

\noindent{\bf Case 1.1}:  \emph{The vertex $u$ is an endpoint of $P$.}


Assume first that $y_1$ is the other endpoint of $P$. Then in $\varphi'$, $u$ is missing $\gamma$, so in particular no $\gamma$-edges appear in $F_1$ (if there were any before, they must have been recolored to $\alpha$). However since these edges all occur after $y_1$ in the order of  $F_1$, and $y_1$ is missing $\alpha$ in $\varphi'$, $F_1$ is still a multifan in $G$ under $\varphi'$. However now Claim~\ref{uwi} is not satisfied for $\varphi'$, contradicting  our choice of $\varphi$.

We may now assume that $y_1$ is not the other endpoint of $P$. Here, $F_1(u, w_0:y_1)$ is a multifan with respect to $\varphi'$ containing no edges of $J_0$. But now the color $\gamma$ is missing at both $u$ and $y_1$ under $\varphi'$, so Claim~\ref{uwi} is not satisfied for $\varphi'$, contradicting our choice of $\varphi$. 


\noindent{\bf Case 1.2}:  \emph{The vertex $u$ is not an endpoint of $P$.}   

In this case we know that $P$ does not contain any edges of $F_1$. 

First assume that $y_1$ is not an endpoint of $P$. Then $F_1$ is still a multifan with respect to $\varphi'$ containing no edges of $J_0$. However the color  $\alpha$ is missing at both $u$ and at least one other vertex from $\{y_2, y_3, y_4\}$ under $\varphi'$. This means that Claim~\ref{uwi} is not satisfied for $\varphi'$, contradicting our choice of $\varphi$.

We may now assume that $y_1$ is an endpoint of $P$. Now, $F_1(u, w_0:y_1)$ is a multifan with respect to $\varphi'$ containing no edges of $J_0$. But then $\alpha$ is missing at both $u$ and $y_1$ under $\varphi'$, so Claim~\ref{uwi} is not satisfied for $\varphi'$, contradicting our choice of $\varphi$.  This completes our proof of (1).

\noindent (2) Suppose for a contradiction that $y_1$ and $y_2$ are two distinct vertices of $\{w_0, \ldots, w_t\}$, both missing $\gamma$ under $\varphi$. Suppose, without loss of generality, that $y_1$ is the first vertex 
in the order $w_0, w_1, \ldots, w_t$ for which $\gamma\in \pbar(y_1)$. In particular, this implies that $F_1(u, w_0:y_1)$ does not contain any edge that is colored with $\gamma$.  

Let $\alpha\in \pbar(u)$; we know that $\alpha\ne \gamma$ by Claim~\ref{uwi}. Consider the set of  maximal alternating $(\alpha, \gamma)$-paths staring at $u, y_1$ or $y_2$. If these three paths are all disjoint, 
then one of them, say $P$,  does not contain any edge of $J_0$ (since there are at most two edges of $J_0$ with the colors $\alpha, \gamma$). On the other hand, suppose that two of the three paths are in fact the same path $P^*$. By our discussion prior to this claim, $P^*$ must contain $x$ as an internal vertex, and hence there are two $J_0$-edges colored $\alpha, \gamma$ on the path $P^*$. But then again, among the three paths, we can find one, say $P$, that does not contain any edges of $J_0$.

We let $\varphi'=\varphi/P$ and consider two different cases according to whether or not $u$ is an endpoint of $P$. Note that in either case, we have ensured that $P$ doesn't contain any edge of $J_0$, does not contain $x$, and only one of its endpoints is from $F_1$.


\noindent{\bf Case 2.1}:  \emph{The vertex $u$ is an endpoint of $P$.}

Since the other endpoint of $P$ is outside $F_1$, $F_1(u, w_0:y_1)$ is a multifan with respect to $\varphi'$ containing no edges of $J_0$. But now the color $\gamma$ is missing at both $u$ and $y_1$ under $\varphi'$, so Claim~\ref{uwi} is not satisfied for $\varphi'$, leading to a contradiction in our choice of $\varphi$.

\noindent{\bf Case 2.2}:  \emph{The vertex $u$ is not an endpoint of $P$.}   
In this case we know that $P$ does not contain any edges of $F_1$, and our proof is identical to that in Case 1.2 above. 
\qed

\vspace*{.2in}

We can now complete the proof of our lemma. For statement (a), Claims~\ref{uwi} and~\ref{gamma} tell us that
\begin{equation}\label{VF1a}
	|\pbar(V(F_1))|  \ge  |\pbar(u)|+ \tfrac{1}{3}\sum_{i=0}^t|\pbar(w_i)|.
\end{equation}
If we let $G'$ be the graph induced by all the colored edges in $G$, then $|\pbar(z)|\geq k-d_{G'}(z)$ for all $z\in V(G)$. Since $d_{G'}(z)\leq \Delta$, we know that $|\pbar(z)|\geq 4$ for any $z$. Since the edge $uw_0$ is uncolored, we further know that $|\pbar(w_0)|\geq 5$ (and so $t\ge 3$). So from (\ref{VF1a}), we get
\begin{equation}\label{pbarLbound}
|\pbar(V(F_1))|  \ge  (k-d_{G'}(u))+ \tfrac{1}{3}(5)+ \tfrac{1}{3}(4t)>k+2+t-d_{G'}(u).
\end{equation}
On the other hand, we know that there are exactly $t$ colored edges in $F_1$,  so there are exactly $d_{G'}(u)-t$ colored edges incident to $u$ that are not included in $F_1$. At most two of these excluded colored edges could be in $J$ (since $u\neq x$) and by the assumption of the lemma, while all the rest,     under $\varphi$,  must be colored by a color not in $\pbar(\{w_0, \ldots, w_t\})$, by (P2) and by definition of $F_1$. Of course, by virtue of being incident to $u$, none of these excluded colors are in $\pbar(u)$ either. Hence we have found a set of at least $d_{G'}(u)-t-2$ colors from $[1,k]$ that are not in $\pbar(V(F_1)$. So
$$|\pbar(V(F_1))| \leq k-(d_{G'}(u)-t-2)=k+2+t-d_{G'}(u),$$
contradicting (\ref{pbarLbound}).

For statement (b),
Claims~\ref{uwi} and~\ref{gamma} tell us that
\begin{equation}\label{VF1}
	|\pbar(V(F_1))|  \ge  |\pbar(u)|+ \sum_{i=0}^t|\pbar(w_i)|. 
\end{equation}
If we let $G'$ be the graph induced by all the colored edges in $G$, then $\pbar(z)\geq k-d_{G'}(z)$ for all $z\in V(G)$. 
We further argued that $|\pbar(w_0)|\geq 2$ earlier in this proof. Since $w_i\in X$ by the assumption that $u\in Y$, we know that $d_{G'}(w_i) \le d_G(w_i) \le k-1$. Thus $|\pbar(w_i)|\geq 1$ for all other $i$. So from (\ref{VF1}), we get
\begin{equation}\label{pbarLbound1}
	|\pbar(V(F_1))|  \ge k-d_{G'}(u)+ 2+ t.
\end{equation}


On the other hand, we know that there are exactly $t$ colored edges in $F_1$,  so there are exactly $d_{G'}(u)-t$ colored edges incident to $u$ that are not included in $F_1$. At most one of these excluded colored edges could be in $J_0$ (since $u\neq x$ and $\Delta(G[J]-x) \le 1$), while all the rest must have colors that are not in $\pbar(\{w_0, \ldots, w_t\})$, by (P2) and by definition of $F_1$. Of course, by virtue of being incident to $u$, none of these excluded colors are in $\pbar(u)$ either. Hence we have found a set of at least $d_{G'}(u)-t-1$ colors from $[1,k]$ that are not in $\pbar(V(F_1))$. So we get 
$$|\pbar(V(F_1))| \leq k-(d_{G'}(u)-t-1)< k-d_{G'}(u)+ 2+ t,$$
contradicting~\eqref{pbarLbound1}. 
\end{proofc}

\end{document}
